\section{Formal Finite Injury}

The finite-injury proof is isolated in \texttt{FiniteInjury.lean}.  Its
goal is not merely to show that records eventually have stable values, but
to provide the exact stabilization package consumed by the requirement
satisfaction proofs.

\subsection{Record Stability and Quietness}

The theorem \texttt{record_step_stable} states that if requirement \(i\)'s
record exists at stage \(s\), and no requirement \(j\leq i\) receives
attention at that stage, then requirement \(i\)'s record is unchanged from
stage \(s\) to stage \(s+1\).  The theorem \texttt{record_stable_of}
extends this to an interval.  This is the no-injury fact in its formal
form: if nothing at or above a requirement acts, its record is frozen.

\subsection{The Rank Argument}

The local progress measure is:

\begin{lstlisting}
def rank (r : ReqState) : Nat :=
  (if r.witness.isSome then 1 else 0) +
  (if r.acted then 1 else 0)

def reqRank (i s : Nat) : Nat :=
  rank (((stageState s).reqs[i]?).getD ReqState.init)
\end{lstlisting}

The theorem \texttt{rank_le_two} bounds every rank by \(2\).
\texttt{reqRank_incr} proves that when requirement \(i\) receives
attention, and its record already exists, its rank increases by exactly
one.  Appointment changes the witness bit from absent to present; action
changes the acted bit from false to true.  The theorem
\texttt{reqRank_mono} proves that if all strictly higher-priority
requirements are quiet at a step, the rank of \(i\) never decreases.

\subsection{Induction Over Priority}

The main theorem is:

\begin{lstlisting}
theorem quiet_above (i : Nat) :
  exists s0,
    forall s, s0 <= s ->
    forall j, j <= i -> attended s != some j
\end{lstlisting}

The proof uses induction on \(i\).  Given a stage after which all
requirements \(j<i\) are quiet, suppose for contradiction that \(i\)
receives attention arbitrarily late.  Choose three later attention stages.
Between them, higher priorities are quiet, so the rank cannot decrease;
at each attention to \(i\), the rank increases.  This forces the rank above
\(2\), contradicting \texttt{reqRank_le_two}.  Hence \(i\) eventually stops
receiving attention, and together with the induction hypothesis all
requirements \(j\leq i\) are eventually quiet.

The theorem \texttt{exists_stable} packages the conclusion: for every
requirement \(i\) there are a stage \(s_1\) and a record \(r\) such that
\(i<s_1\), the record at index \(i\) is always \(r\) after \(s_1\), and no
requirement \(j\leq i\) is ever attended after \(s_1\).  Finally,
\texttt{stable_not_requires} shows that after such a stage the requirement
does not even require attention; otherwise the leastness theorem for
\texttt{attended} would select some \(j\leq i\).

\begin{figure}
\centering
\begin{tikzpicture}[
  node distance=12mm,
  state/.style={draw, rounded corners, align=center, minimum width=32mm, minimum height=8mm, fill=green!6},
  inj/.style={draw, rounded corners, align=center, minimum width=36mm, minimum height=8mm, fill=red!8},
  arr/.style={-Latex, thick}
]
\node[state] (init) {blank\\\texttt{witness = none}};
\node[state, right=of init] (wait) {waiting\\\texttt{some x}, unacted};
\node[state, right=of wait] (acted) {acted\\restrained};
\node[inj, below=of wait] (injury) {higher-priority action\\initializes lower records};

\draw[arr] (init) -- node[above]{appoint fresh witness} (wait);
\draw[arr] (wait) -- node[above]{see output \(0\)} (acted);
\draw[arr] (injury) -- (init);
\draw[arr] (wait) -- (injury);
\draw[arr] (acted) -- (injury);
\end{tikzpicture}


\caption{Lifecycle of one strategy.  Higher-priority action initializes
lower-priority records; once higher priorities are quiet, attention can
only move the strategy forward through a bounded rank.}
\end{figure}

\begin{figure}
\centering
\begin{tikzpicture}[
  node distance=10mm,
  box/.style={draw, rounded corners, align=center, minimum width=42mm, minimum height=9mm, fill=purple!6},
  arr/.style={-Latex, thick}
]
\node[box] (ih) {all \(j<i\) eventually quiet};
\node[box, below=of ih] (rank) {after that stage\\rank of \(i\) cannot decrease};
\node[box, below=of rank] (att) {each attention to \(i\)\\strictly increases rank};
\node[box, below=of att] (bound) {rank \(\le 2\)};
\node[box, below=of bound] (quiet) {therefore \(i\) eventually quiet};
\draw[arr] (ih) -- (rank);
\draw[arr] (rank) -- (att);
\draw[arr] (att) -- (bound);
\draw[arr] (bound) -- (quiet);
\end{tikzpicture}


\caption{Finite-injury induction over priority.}
\end{figure}

\subsection{Satisfaction of Requirements}

\texttt{Requirements.lean} proves \texttt{R_satisfied} and
\texttt{S_satisfied}.  We describe the \(R_e\) case; \(S_e\) is symmetric.
Let \(i=2e\), and take the stable record \(r\) and stage \(s_1\) given by
\texttt{exists_stable}.  The proof first shows that \(r\) must contain a
witness \(x\); otherwise the requirement would require attention by
\texttt{requiresAttention_of_no_witness}, contradicting
\texttt{stable_not_requires}.

If \(r.acted=true\), then the invariant \texttt{acted_in} puts \(x\) in
\(A\), while \texttt{acted_comp} supplies a stored computation of program
\(e\) against the current \(B\)-oracle snapshot with output \(0\) and use
\(r.restraint\).  The theorem \texttt{restraint_respected_B} proves that no
later enumeration below this restraint enters \(B\), using the quiet tail
and the invariant \texttt{wit_ge}.  Then
\texttt{run_halt_limit_of_restraint} transfers the finite computation to
the limiting oracle \(Bset\).  Thus the program outputs \(0\) at \(x\),
while the characteristic function of \(Aset\) outputs \(1\).

If \(r.acted=false\), the invariant \texttt{unacted_out} shows that \(x\)
never enters \(Aset\).  If a candidate reduction computed \(Aset\) from
\(Bset\), it would output \(0\) at \(x\).  By
\texttt{run_halt_snapshot_of_limit}, this convergent computation would be
visible against all sufficiently late finite snapshots.  Choosing a stage
large enough for the fuel and use, the helper
\texttt{run_halt_toFinset_snapshot} converts that visibility into
\texttt{convCheck}, so the requirement would again require attention,
contradicting \texttt{stable_not_requires}.

The final nonreducibility theorems are:

\begin{lstlisting}
theorem not_A_le_B : not (TuringReducible Aset Bset)
theorem not_B_le_A : not (TuringReducible Bset Aset)
\end{lstlisting}

They use the program numbering: an arbitrary candidate program \(c\) is
assigned index \texttt{OracleCode.encodeCode c}, and the round-trip theorem
\texttt{OracleCode.ofNatCode_encodeCode} identifies the corresponding
requirement with \(c\).

